\documentclass[10pt,a4paper]{amsart}
\usepackage[margin=19mm]{geometry}
\usepackage{amsmath,amssymb,mathtools}
\usepackage[scr=rsfs,cal=euler]{mathalfa}
\usepackage{xfrac}
\usepackage[unicode,hidelinks,bookmarks=false]{hyperref}
\newcommand{\ie}{i.e.\ }
\newcommand{\cf}{cf.\ }
\newcommand{\resp}{respectively\ }
\newcommand{\Subsection}{\subsection}
\providecommand{\nobreakdash}{\nobreak\hbox{-}}
\setlength{\emergencystretch}{2em}
\multlinegap0pt
\usepackage{mathtools}
\usepackage{amssymb}
\usepackage{tikz}
\usepackage{enumerate}
\newtheorem{lemma}{Lemma}
\newtheorem{proposition}{Proposition}
\renewcommand{\d}{\ensuremath{\mathrm{d}}}
\title{Strong solutions and stability for a thin-film equation of shear-thinning fluids with contact line in partial wetting}
\author{Manuel V. Gnann, Christina Lienstromberg,, Katerina Nik}
\date{Source: arXiv:2604.22590v1 (CC BY 4.0)}
\begin{document}
\maketitle

\noindent\emph{Source and licence.} Strong solutions and stability for a thin-film equation of shear-thinning fluids with contact line in partial wetting. Version arXiv:2604.22590v1 (CC BY 4.0), distributed by arXiv under the Creative Commons Attribution 4.0 International licence, http://creativecommons.org/licenses/by/4.0/, as recorded in the arXiv licence metadata for this exact version and shown on the abstract page https://arxiv.org/abs/2604.22590v1. Full licence notice: \texttt{licenses/arxiv-2604.22590-CC-BY-4.0.txt}.\par\smallskip

\noindent\textit{Editorial scope.} Counted unit: the article's proposition prop-est-nv (source0.tex lines 2030--2054). The article's complete original proof of it is delivered with it (source0.tex lines 2956--3173, one \texttt{proof} environment of 218 lines, which the article places away from the statement). No mathematics is written, repaired or re-derived: the statement and the proof are the article's own lines, in the article's own notation, and no retained line is substituted or re-flowed.  Retained uncounted as the source-local context the proof needs: lines 306--309; lines 672--674; lines 914--924; lines 919--921; lines 1476--1498; lines 1480--1482; lines 2919--2937.  Nothing was omitted from any retained span, so the package carries no omission record.  No retained line contains a citation, so the wrapper carries no bibliography and no bibliography entry.  No retained line contains a figure environment, an image-inclusion command or drawing code, so no image is delivered and the package declares no asset dependency.  Editorial formatting only, in addition: an AMS \texttt{amsart} preamble that restates the article's own declarations and macros used by the retained lines, namely the environments \texttt{lemma}, \texttt{proposition} on separate counters, together with the standard packages those lines need.  The retained environments print in this document as Lemma 1, Proposition 1 in order of appearance, whereas the article prints the same items with the numbering of its own full document.  The complete native source of the article as distributed in the arXiv e-print for this exact version is bundled unchanged as \texttt{sources/199109/source0.tex}.
\begin{align}
N(u) &\coloneqq (1-  (1+u)^{\frac{3}{2}}) \,  \partial_x^2 (x^{\alpha+2} \, g_{\alpha}(\partial_x^2 u))
+ (1+u)^{\frac{3}{2}}\, \partial_x^2\big( (1- (1+u)^{\frac{\alpha}{2}}) x^{\alpha+2} \, g_{\alpha}(\partial_x^2 u) \big), \label{def-nu}
\end{align}

\begin{lemma}\label{lem-approx-ualpha}
For $\alpha > 1$ the space of test functions $C_\mathrm{c}^\infty((0,\infty))$ is dense in $U_\alpha$.
\end{lemma}

\begin{lemma}\label{lem:weight xu}
For \(\alpha >1\) and \(\beta \in \big[\frac{1-\alpha}{\alpha +1}, \frac{\alpha+2}{5\alpha +3}\big)\),
%
we have
%
\begin{equation}\label{eq:weight xu}
\sup_{x \ge 0} x^\beta |u| \lesssim_{\alpha,\beta} \Big(\int_0^\infty u^2 \, \d x\Big)^{\frac{\alpha -1 + (\alpha +1)\beta}{3\alpha -1}} \Big(\int_0^\infty x^{\alpha+2} \, |\partial_x^2 u|^{\alpha+1} \, \d x\Big)^{\frac{1-2\beta}{3\alpha -1}}
\end{equation}
%
for all $u \in U_\alpha$.
\end{lemma}

\begin{equation}
\sup_{x \ge 0} x^\beta |u| \lesssim_{\alpha,\beta} \Big(\int_0^\infty u^2 \, \d x\Big)^{\frac{\alpha -1 + (\alpha +1)\beta}{3\alpha -1}} \Big(\int_0^\infty x^{\alpha+2} \, |\partial_x^2 u|^{\alpha+1} \, \d x\Big)^{\frac{1-2\beta}{3\alpha -1}}
\end{equation}

\begin{lemma}\label{lem:weight xu lp'}
%
Let $\alpha >1$. Then the following estimate holds:
%
 \begin{equation}\label{eq:weight xu lp'}
\int_0^\infty x^{\nu p} |\partial _x u|^p \, \d x \lesssim_\alpha \Big(\int_0^\infty u^2 \, \d x\Big)^\gamma \Big(\int_0^\infty x^{\alpha+2} \, |\partial_x^2 u|^{\alpha+1} \, \d x\Big)^\delta
 \end{equation}
%
 for all $u \in U_\alpha$, where the exponents $\gamma$ and $\delta$ are given by
 %
 \begin{equation*}
     \gamma = \frac{\alpha+1+((\alpha+1)\nu - 2) p}{3\alpha -1} \quad \text{ and } \quad \delta= \frac{(3-2\nu) p - 2}{3\alpha -1},
 \end{equation*}
 %
where either $p = \frac{4(\alpha+1)}{\alpha+3}$ and $\nu = \frac{\alpha+2}{2(\alpha+1)}$, or $p>\frac{4(\alpha+1)}
%
{\alpha+3}$ and 
%
 \begin{equation*}
\nu \in \Big[\frac{2}{\alpha +1} + \frac{2(\alpha-2)}{(\alpha +3) p}, \frac{\alpha +2}{5\alpha +3} \big(3 - \tfrac 2 p \big)\Big) 
\end{equation*}
%
\end{lemma}

 \begin{equation}
\int_0^\infty x^{\nu p} |\partial _x u|^p \, \d x \lesssim_\alpha \Big(\int_0^\infty u^2 \, \d x\Big)^\gamma \Big(\int_0^\infty x^{\alpha+2} \, |\partial_x^2 u|^{\alpha+1} \, \d x\Big)^\delta
 \end{equation}

\begin{align}\nonumber
    \int_0^\infty N(v)\, u\, \d x
    & \lesssim_\alpha
    B C^\frac{2(\alpha-1)}{3\alpha-1} D^\frac{3\alpha^2+1}{3\alpha-1} 
    +
    A^\frac{\alpha-1}{3\alpha-1}
    B^\frac{2\alpha}{3\alpha-1}
    C^\frac{\alpha-1}{3\alpha-1} 
    D^\frac{\alpha(3\alpha+1)}{3\alpha-1}
    \\
    &\qquad
    + 
    A^\frac{2(\alpha-1)}{3\alpha-1} B^\frac{\alpha+1}{3\alpha-1}
    C^\frac{2(\alpha-1)}{3\alpha-1} 
    D^\frac{3\alpha(\alpha+1)}{3\alpha-1}
    + 
    A^\frac{2(\alpha-1)}{3\alpha-1} B^\frac{\alpha+1}{3\alpha-1} 
    D^{\alpha+1},\label{est-N(v)u}
\end{align}

\begin{proposition}\label{prop-est-nv}
Let $\alpha > 2$ and suppose that $v \in V_\alpha$ with  $\|v\|_{U_\alpha} \ll_\alpha 1$. Then, we have the inequality
%
\begin{align}\label{est-nv}
    \int_0^\infty \big(N(v)\big)^2\, \, \d x
    &\lesssim_\alpha
    A^\frac{2(\alpha-1)}{3\alpha-1}
    B^\frac{2}{3\alpha-1}
    C+
    A^\frac{4(\alpha-1)}{3\alpha-1}
    B^\frac{4}{3\alpha-1}
    C, 
\end{align}
%
where
%
\begin{equation}\label{eq:def_ABC_2}
    A \coloneq 
    \int_0^\infty v^2\, \d x,
    \quad
    B \coloneq \int_0^\infty x^{\alpha+2} |\partial_x^2 v|^{\alpha+1}\, \d x,
    \quad \text{and} \quad
    C \coloneq \int_0^\infty (\partial_x^2 w)^2\, \d x.
\end{equation}
\end{proposition}

\begin{proof}
By approximation according to Lemma~\ref{lem-approx-ualpha}, we may assume $u \in C^\infty_\mathrm{c}((0,\infty))$, thus justifying all integrations by parts within this proof. By \eqref{est-nv} of Proposition~\ref{prop-est-nv} we have $N(v) \in L^2(0,\infty)$, so that the left-hand side of \eqref{est-N(v)u} is well-defined. Using the definition \eqref{def-nu} of $N(\cdot)$ and applying integration by parts twice on each term of $N(v)u$, we obtain
%
\begin{align}\nonumber
    \int_0^{\infty} N(v) u\, \d x 
    &= 
    \int_0^{\infty} u\, \bigl(1-(1+v)^{\frac{3}{2}}\bigr) \, \partial^2_x\bigl(x^{\alpha+2}\, g_{\alpha}(\partial_x^2 v)\bigr) \, \d x 
    \\
    & \quad 
    + \int_0^{\infty} u\,  (1+v)^{\frac{3}{2}} \, \partial^2_x\bigl(
    \bigl(1-(1+v)^{\frac{\alpha}{2}}\bigr) x^{\alpha+2}\, g_{\alpha}(\partial_x^2 v)\bigr) \, \d x 
    \nonumber\\
    &= 
    \int_0^{\infty} x^{\alpha+2}\, g_{\alpha}(\partial_x^2 v)\, 
    \partial^2_x\bigl(u\, \bigl(1-(1+v)^{\frac{3}{2}}\bigr)\bigr) \, \d x 
    \nonumber\\
    &\quad 
    + \int_0^{\infty} \bigl(1-(1+v)^{\frac{\alpha}{2}}\bigr) \, x^{\alpha+2}\, g_{\alpha}(\partial_x^2 v)\, \partial^2_x\bigl(u\,  (1+v)^{\frac{3}{2}}\bigr) \,\d x
    \nonumber\\
    &\eqcolon
    I_1 + I_2. \label{eq:estimate1RHS}
\end{align}
%
We start by estimating $I_1$. Due to estimate~\eqref{eq:weight xu} of Lemma~\ref{lem:weight xu} and $\|v\|_{U_\alpha} \ll_\alpha 1$, we can assume $\sup_{x\geq 0} |v| \le \frac 1 2$. Thus, computing $\partial^2_x\bigl(u\, \bigl(1-(1+v)^{\frac{3}{2}}\bigr)\bigr)$ and using the definition of $g_\alpha$, we may estimate 
%
\begin{align}\nonumber
    I_1
    &=
    \int_0^\infty x^{\alpha+2} g_\alpha(\partial_x^2 v)  \bigl(1 - (1 + v)^\frac{3}{2}\bigr)\, \partial_x^2u\, \d x 
    - 
    3 \int_0^\infty x^{\alpha+2} g_\alpha(\partial_x^2 v) (1 + v)^\frac{1}{2} (\partial_x v)\, (\partial_x u)\, \d x
    \\
    & \phantom{=}
    - \frac{3}{4} \int_0^\infty x^{\alpha+2} g_\alpha(\partial_x^2 v) (\partial_x v)^2\, (1 + v)^{-\frac{1}{2}}\, u\,  \d x
    -
    \frac{3}{2} \int_0^\infty x^{\alpha+2} |\partial_x^2 v|^{\alpha+1} (1 + v)^\frac{1}{2} u\, \d x
    \nonumber\\
    & \lesssim
    \sup_{x\geq 0} |v|
    \int_0^\infty x^{\alpha+2} 
    |\partial_x^2 v|^\alpha\, |\partial_x^2 u|\, \d x 
    + 
    \int_0^\infty x^{\alpha+2} 
    |\partial_x^2 v|^\alpha\, |\partial_x v|\, |\partial_x u|\, \d x 
    \nonumber\\
    & \phantom{=} 
    +
    \int_0^\infty x^{\alpha+2} |\partial_x^2 v|^\alpha\, (\partial_x v)^2\, |u|\, \d x
    + 
    \int_0^\infty x^{\alpha+2} |\partial_x^2 v|^{\alpha+1}\, |u|\, \d x
    \nonumber\\
    &=
    J_1 + J_2 + J_3 + J_4. \label{eq:N(v)u_J_i}
\end{align}
%
Applying H\"older's inequality to $J_1$ and estimating $\sup_{x\geq 0} |v|$ by means of Lemma~\ref{lem:weight xu} with $\beta=0$, we find
%
\begin{align}\nonumber
    J_1 
    &\leq_{\phantom{\alpha}} 
    \Big(\int_0^\infty x^{\alpha+2} |\partial_x^2 u|^{\alpha+1}\, \d x\Big)^\frac{1}{\alpha+1} \sup_{x \geq 0} |v|\, \Big(\int_0^\infty x^{\alpha+2} |\partial_x^2 v|^{\alpha+1}\, \d x\Big)^\frac{\alpha}{\alpha+1}
    \\
    &\stackrel{\mathclap{\eqref{eq:weight xu}}}{\lesssim}_\alpha 
    \Big(\int_0^\infty x^{\alpha+2} |\partial_x^2 u|^{\alpha+1}\, \d x\Big)^\frac{1}{\alpha+1}
    \Big(\int_0^\infty v^2\, \d x\Big)^\frac{\alpha-1}{3\alpha-1}
    \Big(\int_0^\infty x^{\alpha+2} |\partial_x^2 v|^{\alpha+1}\, \d x\Big)^\frac{3\alpha^2+1}{(3\alpha-1)(\alpha+1)}. \label{eq:N(v)u_J1}
\end{align}
%
In order to estimate the integral $J_2$, we first apply H\"older's inequality which gives
%
\begin{equation*}
    J_2
    \leq
    \Big(\int_0^\infty x^{\alpha+2} |\partial_x u|^{2(\alpha+1)}\, \d x\Big)^\frac{1}{2(\alpha+1)}
    \Big(\int_0^\infty x^{\alpha+2} |\partial_x^2 v|^{\alpha+1}\, \d x\Big)^\frac{\alpha}{\alpha+1} 
    \Big(\int_0^\infty x^{\alpha+2} |\partial_x v|^{2(\alpha+1)}\, \d x\Big)^\frac{1}{2(\alpha+1)}.
\end{equation*}
%
The second factor is benign. To estimate the other two factors, we apply Lemma~\ref{lem:weight xu lp'} with $p = 2(\alpha+1) > \frac{4(\alpha+1)}{\alpha+3}$ and $\nu = \frac{\alpha+2}{2(\alpha+1)}$, where we note that
%
\begin{align*}
\nu &\ge \frac{2}{\alpha+1} + \frac{2(\alpha-2)}{(\alpha+3)p} = \frac{3\alpha+4}{(\alpha+1)(\alpha+3)} \quad \Longleftarrow \quad \alpha \ge 2, \\
\nu &< \frac{\alpha+2}{5\alpha+3} \big(3-\tfrac 2 p\big) = \frac{(\alpha+2)(3\alpha+2)}{(\alpha+1)(5\alpha+3)} \quad \Longleftarrow \quad \alpha > -1.
\end{align*}
%
This gives
%
\begin{equation*}
    \int_0^\infty x^{\alpha+2} |\partial_x u|^{2(\alpha+1)}\, \d x
    \; \stackrel{\mathclap{\eqref{eq:weight xu lp'}}}{\lesssim}_\alpha \;
    \Big(\int_0^\infty u^2\, \d x\Big)^\frac{(\alpha-1)(\alpha+1)}{3\alpha-1}
    \Big(
    \int_0^\infty x^{\alpha+2} |\partial_x^2 u|^{\alpha+1}\, \d x
    \Big)^\frac{4\alpha}{3\alpha-1}
\end{equation*}
%
for the first factor and the same, with $u$ replaced by $v$, for the third factor. Thus, for $J_2$ we end up with
\begin{align}\nonumber
    J_2
    & \lesssim_\alpha
    \Big(\int_0^\infty
    u^2\, \d x
    \Big)^\frac{\alpha-1}{2(3\alpha-1)}
    \Big(\int_0^\infty
    x^{\alpha+2} |\partial_x^2 u|^{\alpha+1}\, \d x
    \Big)^\frac{2\alpha}{(\alpha+1)(3\alpha-1)} \\
    & \phantom{\lesssim_\alpha} \times \Big(\int_0^\infty
    v^2\, \d x
    \Big)^\frac{\alpha-1}{2(3\alpha-1)} 
    \Big(\int_0^\infty
    x^{\alpha+2} |\partial_x^2 v|^{\alpha+1}\, \d x
    \Big)^\frac{\alpha(3\alpha+1)}{(\alpha+1)(3\alpha-1)}. \label{eq:N(v)u_J2}
\end{align}
%
Next, we estimate $J_3$. We have
%
\begin{equation*}
\begin{split}
    J_3
    \leq
    \sup_{x \geq 0} |u| 
    \int_0^\infty x^{\alpha+2} |\partial_x^2 v|^\alpha\, (\partial_x v)^2\, \d x
\end{split}
\end{equation*}
%
and observe that $\sup_{x \geq 0} |u|$ can again be estimated by estimate~\eqref{eq:weight xu} of Lemma~\ref{lem:weight xu} with $\beta=0$, while the integral can be treated like $J_2$ with $u=v$. This leads to the estimate
%
\begin{equation}\label{eq:N(v)u_J3}
    J_3
    \lesssim_\alpha
    \Big(\int_0^\infty u^2\, \d x\Big)^\frac{\alpha-1}{3\alpha-1}
    \Big(\int_0^\infty x^{\alpha+2} |\partial_x^2 u|^{\alpha+1} \, \d x\Big)^\frac{1}{3\alpha-1}
    \Big(\int_0^\infty
    v^2\, \d x
    \Big)^\frac{\alpha-1}{3\alpha-1} 
    \Big(\int_0^\infty
    x^{\alpha+2} |\partial_x^2 v|^{\alpha+1}\, \d x
    \Big)^\frac{3\alpha}{3\alpha-1}.
\end{equation}
%
By applying Lemma~\ref{lem:weight xu} once more, we obtain for $J_4$,
%
\begin{equation} \label{eq:N(v)u_J4}
    J_4 
    \, \stackrel{\mathclap{\eqref{eq:weight xu}}}{\lesssim}_\alpha 
    \Big(\int_0^\infty u^2\, \d x\Big)^\frac{\alpha-1}{3\alpha-1}
    \Big(\int_0^\infty x^{\alpha+2} |\partial_x^2 u|^{\alpha+1} \, \d x\Big)^\frac{1}{3\alpha-1}
    \int_0^\infty x^{\alpha+2} |\partial_x^2 v|^{\alpha+1}\, \d x.
\end{equation}
%
Finally, using \eqref{eq:N(v)u_J1}--\eqref{eq:N(v)u_J4} in \eqref{eq:N(v)u_J_i} we end up with
%
\begin{align}\nonumber
    I_1
    &\lesssim_\alpha
    \Big(\int_0^\infty x^{\alpha+2} |\partial_x^2 u|^{\alpha+1}\, \d x\Big)^\frac{1}{\alpha+1}
    \Big(\int_0^\infty v^2\, \d x\Big)^\frac{\alpha-1}{3\alpha-1}
    \Big(\int_0^\infty x^{\alpha+2} |\partial_x^2 v|^{\alpha+1}\, \d x\Big)^\frac{3\alpha^2+1}{(3\alpha-1)(\alpha+1)}
    \\
    &\phantom{\lesssim_\alpha}
    +
    \Big(\int_0^\infty
    u^2\, \d x
    \Big)^\frac{\alpha-1}{2(3\alpha-1)}
    \Big(\int_0^\infty
    x^{\alpha+2} |\partial_x^2 u|^{\alpha+1}\, \d x
    \Big)^\frac{2\alpha}{(\alpha+1)(3\alpha-1)} \nonumber \\
    &\phantom{\lesssim_\alpha +} \times \Big(\int_0^\infty
    v^2\, \d x
    \Big)^\frac{\alpha-1}{2(3\alpha-1)} 
    \Big(\int_0^\infty
    x^{\alpha+2} |\partial_x^2 v|^{\alpha+1}\, \d x
    \Big)^\frac{\alpha(3\alpha+1)}{(\alpha+1)(3\alpha-1)}
    \nonumber \\
    &\phantom{\lesssim_\alpha}
    +
    \Big(\int_0^\infty u^2\, \d x\Big)^\frac{\alpha-1}{3\alpha-1}
    \Big(\int_0^\infty x^{\alpha+2} |\partial_x^2 u|^{\alpha+1} \, \d x\Big)^\frac{1}{3\alpha-1}
    \Big(\int_0^\infty
    v^2\, \d x
    \Big)^\frac{\alpha-1}{3\alpha-1} 
    \Big(\int_0^\infty
    x^{\alpha+2} |\partial_x^2 v|^{\alpha+1}\, \d x
    \Big)^\frac{3\alpha}{3\alpha-1}
    \nonumber \\
    &\phantom{\lesssim_\alpha}
    +
    \Big(\int_0^\infty u^2\, \d x\Big)^\frac{\alpha-1}{3\alpha-1}
    \Big(\int_0^\infty x^{\alpha+2} |\partial_x^2 u|^{\alpha+1} \, \d x\Big)^\frac{1}{3\alpha-1}
    \int_0^\infty x^{\alpha+2} |\partial_x^2 v|^{\alpha+1}\, \d x. \label{eq:N(v)u_I1}
\end{align} 
%
We obtain a similar estimate for $I_2$. Indeed, recalling the definition of $g_\alpha$ and the assumption $\sup_{x\geq 0} |v| \leq \frac{1}{2} < 1$ because of \eqref{eq:weight xu} of Lemma~\ref{lem:weight xu} and $\|v\|_{U_\alpha} \ll_\alpha 1$, we  apply the mean-value theorem to the first factor and calculate $\partial^2_x\bigl(u\,  (1+v)^{\frac{3}{2}}\bigr)$. This yields
%
\begin{equation*}
\begin{split}
    I_2
    &=
    \int_0^{\infty} \bigl(1-(1+v)^{\frac{\alpha}{2}}\bigr) \, x^{\alpha+2}\, g_{\alpha}(\partial_x^2 v)\, \partial^2_x\bigl(u\,  (1+v)^{\frac{3}{2}}\bigr) \,\d x
    \\
    &\lesssim_\alpha
    \sup_{x \geq 0} |v|
    \Big(
    \int_0^\infty x^{\alpha+2} |\partial_x^2 v|^\alpha\, |\partial_x^2u|\, \d x 
    + 
    \int_0^\infty x^{\alpha+2} |\partial_x^2 v|^\alpha |\partial_x v|\, |\partial_x u|\, \d x
    \\
    &\phantom{\lesssim_\alpha
    \sup_{x \geq 0} |v| \Big(}
    + \int_0^\infty x^{\alpha+2} |\partial_x^2 v|^\alpha (\partial_x v)^2\, |u|\, \d x
    +
    \int_0^\infty x^{\alpha+2} |\partial_x^2 v|^{\alpha+1}\, |u|\, \d x
    \Big).
\end{split}
\end{equation*}
%
The second and third line contain exactly the same integrals $J_1$, $J_2$, $J_3$, and $J_4$ as in estimate~\eqref{eq:N(v)u_J_i} for $I_1$, except the additional pre-factor $\sup_{x \geq 0} |v|$ that appears now in front of all integrals and not only in the first one. However, since $\sup_{x \ge 0} |v| \le \frac 1 2$, estimate~\eqref{eq:N(v)u_I1} is also valid for $I_2$ instead of $I_1$, thus finishing the proof of \eqref{est-N(v)u}.
\end{proof}

\end{document}
